\chapter{Chiral Nucleation and Baryon Asymmetry in a Torsional Substrate}
\label{app:baryon_asymmetry}
\addtocontents{toc}{\protect\setcounter{tocdepth}{0}}

\noindent\textbf{Abstract:} The Standard Model fails to explain the predominance of matter over antimatter (Baryon Asymmetry), as perfectly symmetric quantum models predict equal amounts should have formed and annihilated. Tension Field Theory (TFT) resolves this by identifying the Spacetime Gravitational Floor (SGF) as a Chiral Substrate. Because the universe decompressed from a rotating Black Hole seed, the conservation of angular momentum imparted a global Torsion Field to the vacuum. Consequently, the nucleation energy for matter (aligned with the structural twist) is systematically lower than for antimatter (anti-aligned). This mechanical symmetry-breaking allowed a residual population of matter to survive the initial decompression epoch.

\section{The Physics of the Seed (Origin of Torsion)}
In standard cosmology, the Big Bang is modeled as a formless, isotropic explosion. In TFT, the genesis event is the Catastrophic Decompression of a saturated State II core (a Black Hole).

Because macroscopic black holes inherently possess immense angular momentum (Kerr black holes), this genesis seed was rotating.
\begin{itemize}
    \item \textbf{Conservation of Spin:} The angular momentum of the parent seed is conserved directly within the lattice structure of the offspring universe.
    \item \textbf{The Torsional Lattice:} As the zero-void Simple Cubic core snapped back into the Face-Centered Cubic (FCC) active vacuum, this spin was imparted as a global Torsion Field. The resulting SGF lattice is not a perfectly neutral geometric grid; it is structurally twisted, establishing a ``Resting State'' with a preferred macroscopic helicity (left-handed grain).
\end{itemize}

\section{Topological Formation Energy (The Chiral Knot)}
As established in Section 3, fundamental particles are topological phase defects (chiral knots) formed during the shattering of the lattice. Matter and antimatter are exact topological mirror images.

In a perfectly neutral, untwisted vacuum, their base formation energy ($E_{\mathrm{base}}$) would be identical. However, in a torsional vacuum, the lattice impedance diverges depending on the direction of the knot's twist.

The equations of formation are therefore strictly asymmetrical:
\begin{equation}
    E_{\mathrm{matter}} = E_{\mathrm{base}} - E_{\mathrm{torsion}}
\end{equation}
\begin{equation}
    E_{\mathrm{antimatter}} = E_{\mathrm{base}} + E_{\mathrm{torsion}}
\end{equation}

\begin{itemize}
    \item \textbf{Matter (The Survivor):} Forms with the helical grain of the universe. The inherent lattice twist mechanically assists the knot formation, resulting in a lower nucleation energy state.
    \item \textbf{Antimatter (The Victim):} Forms against the helical grain. The lattice actively resists this inverse knot formation, resulting in a higher nucleation energy state.
\end{itemize}

\section{The Statistical Survivor (Thermodynamic Bias)}
During the extreme high-temperature phase of Catastrophic Decompression, the creation of these topological defects was governed by statistical thermodynamics. In a chaotic thermal environment, the statistical abundance ($N$) of any structural state forming is exponentially suppressed by its mechanical energy barrier, as defined by the Boltzmann distribution:
\begin{equation}
    N \propto e^{-E / k_{\mathrm{B}} T}
\end{equation}

Because the energy barrier for matter was systematically lower ($E_{\mathrm{matter}} < E_{\mathrm{antimatter}}$), the decompression naturally produced a slight statistical excess of matter. The survivor ratio evaluates as:
\begin{equation}
    \frac{N_{\mathrm{matter}}}{N_{\mathrm{antimatter}}} = \frac{e^{-(E_{\mathrm{base}} - E_{\mathrm{torsion}}) / k_{\mathrm{B}} T}}{e^{-(E_{\mathrm{base}} + E_{\mathrm{torsion}}) / k_{\mathrm{B}} T}} = \exp\left( \frac{2E_{\mathrm{torsion}}}{k_{\mathrm{B}} T} \right)
\end{equation}

Even an infinitesimally small torsion energy ($E_{\mathrm{torsion}}$) relative to the massive decompression temperature ($T$) naturally yields the exact one-part-per-billion excess of matter that we observe today.

\section{Conclusion}
The universe is not perfectly symmetric; it is structurally handed. We live in a matter-dominant universe precisely because we were born from a spinning seed. Antimatter is not mysteriously ``missing'' from the cosmos; it was simply mechanically disadvantaged during the initial crystallization of reality.
